%% kapa: exam=1 round=37 year=2026 period=1 form=A subject=real_estate subject_ko=부동산학원론 seq=21 number=101
%% source: Q-Net 감정평가사 기출문제 (cst00309) artlSeq=5258924 file=raw/1008/5258924/2026년 제37회 감정평가사 1차 1교시.pdf
%% conversion: kordoc 4.24.51 -> tools/gen1.py (auto); review=proofread-sonnet; flags=none
\begin{question}{101}
투자규모와 투자기간이 동일한 상호배타적인 부동산 투자안 A, B가 있다. 투자시점에만 현금유출이 발생하고 보유기간 및 매각시점에는 현금유입만 발생한다. 각 투자안은 보유기간 동안 순영업소득(NOI)이 일정하다. A의 내부수익률(IRR)은 12 \%, B의 내부수익률은 15 \%이다. 투자자의 요구수익률이 8 \%일 때, 순현재가치(NPV)는 A가 B보다 크다. 이에 관한 설명으로 옳은 것은? (단, 주어진 조건에 한함)
\begin{choices}
  \item 투자자의 요구수익률이 12 \%일 때, B의 순현재가치는 ‘0’이다.
  \item 투자자의 요구수익률이 15 \%일 때, 순현재가치는 A가 B보다 크다.
  \item 투자자의 요구수익률이 8 \%보다 클 때, 순현재가치는 A가 B보다 항상 크다.
  \item 투자자의 요구수익률이 12 \% 초과 15 \% 미만의 구간에서 순현재가치는 B가 A보다 크다.
  \item 두 투자안의 순현재가치를 같게 하는 투자자의 요구수익률은 12 \% 초과 15 \% 미만의 구간에 위치한다.
\end{choices}
\end{question}
\answer{4}
